\subsection{Genetic Algorithm for Feature Selection}

The genetic algorithm (GA) was employed on the photovoltaic band-gap task to identify an
optimal subset of composition descriptors while maximizing predictive performance as measured
by ROC-AUC. The GA was configured with the following parameters: population size of 100
individuals, 20 generations, crossover probability of 0.8, mutation probability of 0.2, an
elite size of 5 and a feature-count penalty $\lambda = 0.05$. Fitness was evaluated with 3-fold
cross-validation on a stratified sample of 6,000 training compositions; seven model families
were compared first, and the best one (\GapBestModel{}) was used for the final model.

\begin{table}[H]
\centering\small
\begin{tabular}{lcccc}
\toprule
Model & ROC-AUC & PR-AUC & F1 & Training time\\
\midrule
RF & 0.9073 & 0.4663 & 0.2967 & 27.2\,s \\
XGBoost & 0.8997 & 0.4194 & 0.2226 & 7.7\,s \\
ExtraTrees & 0.8968 & 0.4465 & 0.3849 & 24.3\,s \\
LightGBM & 0.8953 & 0.4008 & 0.3793 & 6.9\,s \\
kNN & 0.8435 & 0.3418 & 0.1910 & 3.7\,s \\
SVM & 0.8135 & 0.1999 & 0.0000 & 86.6\,s \\
LogReg & 0.7986 & 0.1801 & 0.2361 & 5.0\,s \\
\bottomrule
\end{tabular}
\caption*{Model comparison on all 142 descriptors (photovoltaic band gap, test set).}
\end{table}


\subsubsection{Evolution of Fitness and Feature Count}

The GA demonstrated effective convergence over 20 generations (Figure~\ref{fig:ga}). The
fitness of the best individual increased steadily from \GAInitFit{} to \GAFinalFit{}, and the
cross-validated ROC-AUC of the best descriptor subset from \GAInitAUC{} to \GAFinalAUC{}
(+2.0\%), above the \GAAllCV{} obtained with all 142 descriptors under the same protocol.
Concurrently, the average number of selected descriptors in the population decreased from
\GAInitAvgN{} to \GAFinalAvgN{}, and the final individual keeps only \GAFinalN{} of the 142
descriptors, a 63.4\% reduction in feature dimensionality.

\begin{figure}[H]
  \centering
  \safeimg{\textwidth}{results/band_gap/band_gap_ga_evolution.png}
  \caption{Evolution of fitness (ROC-AUC) and feature count during genetic algorithm
  optimization. (A) Fitness values showing steady improvement over generations. (B) Feature
  count evolution demonstrating effective feature reduction while maintaining predictive
  performance.}\label{fig:ga}
\end{figure}

\subsubsection{Feature Selection Performance}

The GA successfully identified a compact yet informative descriptor subset. On the 15,633
held-out test compositions, the final model reaches an ROC-AUC of \GapGAAUC{} with only
\GAFinalN{} descriptors, compared with \GapAllAUC{} using all 142 descriptors: a favorable
trade-off between model complexity and predictive performance. Because positives are rare
(\GapPos{}\%), the decision threshold was then calibrated by maximizing F1 on out-of-fold
predictions of the \emph{training} set ($t=\GapThr{}$), which raised the test F1 from
\GapGAFone{} to \GapTunedFone{} and the recall from \GapGARec{} to \GapTunedRec{} without
changing the ROC-AUC.

\begin{table}[H]
\centering\small
\begin{tabular}{lccc}
\toprule
Metric & Initial & Final & Change\\
\midrule
Best fitness & \GAInitFit{} & \GAFinalFit{} & +2.4\%\\
CV ROC-AUC of best subset & \GAInitAUC{} & \GAFinalAUC{} & +2.0\%\\
Number of features (all $\rightarrow$ GA) & 142 & \GAFinalN{} & $-$63.4\%\\
Feature Count (population average) & \GAInitAvgN{} & \GAFinalAvgN{} & $-$27.3\%\\
Minimum Fitness & \GAInitMinFit{} & \GAFinalMinFit{} & +2.8\%\\
\midrule
Test ROC-AUC (all vs GA descriptors) & \GapAllAUC{} & \GapGAAUC{} & $-$0.4\%\\
Test F1 (GA, threshold 0.5 vs calibrated) & \GapGAFone{} & \GapTunedFone{} & +50.1\%\\
\bottomrule
\end{tabular}
\caption*{Genetic Algorithm Feature Selection Results}
\end{table}

\subsubsection{Descriptor Category Analysis}

The selected \GAFinalN{} descriptors represent diverse elemental properties:
\begin{itemize}
  \item \textbf{Valence electrons} (23 descriptors)
  \item \textbf{Periodic position} (11 descriptors)
  \item \textbf{Thermal / mass} (5 descriptors)
  \item \textbf{Elemental band gap / magnetism} (5 descriptors)
  \item \textbf{Electronegativity} (4 descriptors)
  \item \textbf{Radius / volume} (2 descriptors)
  \item \textbf{Stoichiometry} (2 descriptors)
  \item \textbf{Structure (space group)} (1 descriptors)
\end{itemize}

The dominance of valence-electron descriptors (number of filled and unfilled $s$, $p$, $d$, $f$
states) is physically sound: the band gap is set by the occupation of the valence and
conduction bands. The feature reduction from 142 to \GAFinalN{} descriptors improves model
interpretability and reduces computational requirements while maintaining competitive
predictive performance.

\subsection{Applicability Domain Assessment}

To ensure reliable predictions and assess model generalizability, multiple applicability
domain (AD) methods were employed in the space of the \GAFinalN{} GA-selected descriptors. The
AD defines the chemical space where the model provides trustworthy predictions, based on the
training data characteristics.

\subsubsection{PCA-Based Applicability Domain}
Principal Component Analysis (PCA) was performed to visualize the chemical-space distribution
(Figure~\ref{fig:adpca}). The first two principal components explained \ADPCAVar{}\% of total
variance (PC1: \ADPCone{}\%, PC2: \ADPCtwo{}\%). While this moderate variance explanation
reflects the high dimensionality of the descriptor space, the projection reveals consistent
distributions between training and test sets.

\begin{figure}[H]
  \centering
  \safeimg{\textwidth}{results/band_gap/band_gap_ad_pca.png}
  \caption{PCA-based applicability domain analysis. (A) Training set distribution in PCA
  space. (B) Test set projected onto the same PCA space, showing good overlap with training
  distribution.}\label{fig:adpca}
\end{figure}

\subsubsection{K-Means Clustering Analysis}
K-means clustering with $k=6$ was applied to partition the chemical space
(Figure~\ref{fig:adkm}). Test compositions were assigned to the nearest training cluster
centroids, providing a cluster-based AD assessment. The cluster proportions are almost
identical for training and test sets (for example 29.3\% vs 28.8\% for the largest cluster
and 8.8\% vs 8.9\% for the smallest), indicating consistent chemical-space coverage.

\begin{figure}[H]
  \centering
  \safeimg{\textwidth}{results/band_gap/band_gap_ad_kmeans.png}
  \caption{K-means clustering for applicability domain assessment. (A) Training set
  clustering with 6 clusters. (B) Test set assigned to training clusters, showing consistent
  chemical space coverage.}\label{fig:adkm}
\end{figure}

\subsubsection{Distance-Based AD Classification}
A distance-based approach was implemented using the 95th percentile of the training
nearest-neighbour distances as the AD threshold ($D \le \ADThr{}$, Figure~\ref{fig:addist}).
Of the \ADNTest{} test compositions, \ADNIn{} (\ADFracIn{}\%) fell within the applicability
domain, while \ADNOut{} compositions (\ADFracOut{}\%) were identified as outliers.

\begin{figure}[H]
  \centering
  \safeimg{\textwidth}{results/band_gap/band_gap_ad_distance.png}
  \caption{Distance-based applicability domain classification. (A) Distance distribution
  comparison between training and test sets. (B) AD classification showing \ADFracIn{}\% of
  test compositions within the applicability domain.}\label{fig:addist}
\end{figure}

\subsubsection{AD Concordance and Reliability}
The multi-method AD approach demonstrated strong concordance:
\begin{itemize}
  \item \textbf{High Coverage:} \ADFracIn{}\% of test compositions within the AD.
  \item \textbf{Conservative Threshold:} the 95th percentile ensures reliable predictions.
  \item \textbf{Consistent Results:} all three methods (PCA, clustering, distance-based)
  indicated similar AD boundaries.
  \item \textbf{Reliability check:} inside the AD the model retrieves \ADRecInT{} of the
  photovoltaic compositions (F1 = \ADFoneInT{}), outside only \ADRecOutT{} (F1 =
  \ADFoneOutT{}); the precision--recall AUC also drops from \ADPRIn{} to \ADPROut{}.
\end{itemize}

Compositions identified outside the AD represent chemistries distinct from the training space
and should be treated with caution; the web application therefore reports the AD status with
every prediction.

\clearpage
\subsection{Unsupervised Analysis: Association Rules and Clustering}

\subsubsection{Association Rules}
Apriori mining found \RulesStab{} rules for stability, \RulesMag{} for magnetism, \RulesMetal{}
for metallicity and \RulesGap{} for the photovoltaic band gap (minimum support 0.2--0.3\%,
lift $\ge 1.2$, confidence $\ge \min(0.5, 2\times\text{base rate})$). Without any
supervised model, the rules rediscover textbook chemistry:

\begin{table}[H]
\centering\scriptsize
\begin{subtable}{0.49\textwidth}\centering\begin{tabular}{lccc}
\toprule
Rule & Support & Confidence & Lift\\
\midrule
\{F, Mn\} $\Rightarrow$ MAGNETIC & 0.0035 & 0.954 & 4.64 \\
\{Mn, O\} $\Rightarrow$ MAGNETIC & 0.0193 & 0.951 & 4.62 \\
\{Li, Mn\} $\Rightarrow$ MAGNETIC & 0.0072 & 0.949 & 4.61 \\
\{Fe, Li\} $\Rightarrow$ MAGNETIC & 0.0063 & 0.938 & 4.56 \\
\{Cr, Li\} $\Rightarrow$ MAGNETIC & 0.0037 & 0.922 & 4.48 \\
\{Fe, O\} $\Rightarrow$ MAGNETIC & 0.0167 & 0.906 & 4.40 \\
\bottomrule
\end{tabular}

\caption{Magnetism (base rate \MagPos{}\%)}\end{subtable}\hfill
\begin{subtable}{0.49\textwidth}\centering\begin{tabular}{lccc}
\toprule
Rule & Support & Confidence & Lift\\
\midrule
\{Er, Y\} $\Rightarrow$ STABLE & 0.0030 & 0.658 & 2.23 \\
\{Dy, Y\} $\Rightarrow$ STABLE & 0.0031 & 0.648 & 2.20 \\
\{La, Pr\} $\Rightarrow$ STABLE & 0.0031 & 0.640 & 2.17 \\
\{Sm, Y\} $\Rightarrow$ STABLE & 0.0030 & 0.635 & 2.16 \\
\{La, Nd\} $\Rightarrow$ STABLE & 0.0030 & 0.635 & 2.16 \\
\{Tb, Y\} $\Rightarrow$ STABLE & 0.0032 & 0.633 & 2.15 \\
\bottomrule
\end{tabular}

\caption{Stability (base rate \StabPos{}\%)}\end{subtable}\\[0.6em]
\begin{subtable}{0.49\textwidth}\centering\begin{tabular}{lccc}
\toprule
Rule & Support & Confidence & Lift\\
\midrule
\{Pm\} $\Rightarrow$ METAL & 0.0051 & 0.968 & 2.03 \\
\{Ga, Ni\} $\Rightarrow$ METAL & 0.0032 & 0.933 & 1.96 \\
\{Ir\} $\Rightarrow$ METAL & 0.0166 & 0.894 & 1.87 \\
\{Al, Ni\} $\Rightarrow$ METAL & 0.0032 & 0.887 & 1.86 \\
\{Rh\} $\Rightarrow$ METAL & 0.0213 & 0.883 & 1.85 \\
\{Pd\} $\Rightarrow$ METAL & 0.0233 & 0.833 & 1.75 \\
\bottomrule
\end{tabular}

\caption{Metallicity (base rate \MetalPos{}\%)}\end{subtable}\hfill
\begin{subtable}{0.49\textwidth}\centering\begin{tabular}{lccc}
\toprule
Rule & Support & Confidence & Lift\\
\midrule
\{Se\} $\Rightarrow$ PV\_GAP & 0.0087 & 0.179 & 2.69 \\
\{O, V\} $\Rightarrow$ PV\_GAP & 0.0060 & 0.163 & 2.45 \\
\{Fe, O\} $\Rightarrow$ PV\_GAP & 0.0065 & 0.144 & 2.17 \\
\{Br\} $\Rightarrow$ PV\_GAP & 0.0035 & 0.141 & 2.12 \\
\{S\} $\Rightarrow$ PV\_GAP & 0.0113 & 0.135 & 2.03 \\
\{V\} $\Rightarrow$ PV\_GAP & 0.0067 & 0.134 & 2.02 \\
\bottomrule
\end{tabular}

\caption{PV band gap (base rate \GapPos{}\%)}\end{subtable}
\caption*{Top association rules (element sets $\Rightarrow$ property).}
\end{table}

\begin{itemize}
  \item \textbf{Magnetism:} Mn, Fe, Cr and Co compounds, especially oxides and fluorides, are
  magnetic in more than 85\% of the cases (lift $\approx$ 4--4.6), reflecting their partially
  filled $3d$ shells.
  \item \textbf{Stability:} pairs of rare-earth elements (Er--Y, Dy--Y, La--Pr, Nd--Sm) are
  about 2.2 times more often stable than average: these chemically similar elements substitute
  easily for each other in stable structures.
  \item \textbf{Metallicity:} compounds of platinum-group metals (Ir, Rh, Pd, Os, Ru) are metals
  in 83--89\% of the cases.
  \item \textbf{Photovoltaic band gap:} selenides, sulfides, bromides and vanadium/iron oxides
  are enriched 2.0--2.7 times in the photovoltaic window, the classic absorber chemistries.
\end{itemize}

On GNoME alone (no property), element co-occurrence rules such as Zr $\Leftrightarrow$ Hf (lift
4.9), Co $\Leftrightarrow$ P (2.9) and Pr $\Leftrightarrow$ Nd (2.5) reveal the chemical
families explored by the generative search.

\begin{figure}[H]
  \centering
  \safeimg{0.75\textwidth}{results/unsupervised/association_rules_magnetism.png}
  \caption{Top association rules for magnetism (bar length = lift, colour = confidence).}
\end{figure}

\subsubsection{Clustering of the Materials Space}
PCA followed by K-Means was applied to 100,000 GNoME compositions and 25,000 compositions of
each labeled dataset (175,000 points). The silhouette criterion selects $k=\KBest{}$
(Figure~\ref{fig:km}), and the three clusters, found without any label, correspond to
well-known material families: (i) rare-earth / noble-metal intermetallics (92\% metals),
(ii) $p$-block and transition-metal intermetallics (85\% metals), and (iii) ionic compounds
dominated by O, F, S (only 22\% metals but 31\% magnetic and 9.6\% in the photovoltaic window).
The unsupervised structure is therefore strongly aligned with the target properties, which is
the \emph{cluster assumption} exploited by Cluster-then-Label and by the graph-based methods.

\begin{figure}[H]
  \centering
  \safeimg{\textwidth}{results/unsupervised/kmeans_pca_clusters.png}
  \caption{K-Means clusters in PCA space (left) and position of the labeled datasets with
  respect to the unlabeled GNoME pool (right).}\label{fig:km}
\end{figure}

\begin{figure}[H]
  \centering
  \safeimg{0.85\textwidth}{results/unsupervised/kmeans_cluster_property_rates.png}
  \caption{Property rates inside each K-Means cluster.}
\end{figure}

\subsection{Self-Supervised Pre-training}
The masked autoencoder was trained for 20 epochs on 526,352 GNoME compositions (5\% held out).
It reconstructs the full descriptor vector with a validation $R^2$ of \AERtwo{} and the
\emph{masked} (hidden) descriptors with $R^2 = \AERtwoMasked{}$, showing that the 32-dimensional
embedding captures the correlations between elemental properties. Its downstream value is
reported in the result tables as the ``Self-Supervised (AE)'' rows: concatenating the embedding
to the raw descriptors keeps the performance of the tuned baseline, whereas the embedding alone
loses information compared with the 142 descriptors.

\begin{figure}[H]
  \centering
  \safeimg{0.8\textwidth}{unsupervised/autoencoder_training.png}
  \caption{Self-supervised masked-autoencoder training (reconstruction loss and validation
  $R^2$).}
\end{figure}
